\chapter{行动、计划与直接效应}
% 第四章 4.1—4.3（其余见 ch04b.tex）

\begin{flushright}
\kaishu\small
行动超出其智慧者，其智慧必将长存。\\[2pt]
——拉比哈尼拿·本·多萨（Rabbi Hanina ben Dosa，公元 1 世纪）
\end{flushright}

\section*{序言}
\addcontentsline{toc}{section}{序言}

迄今为止，我们对因果效应的分析一直聚焦于 $\doop(x)$ 形式的原始干预——它表示把变量 $X$ 的取值设定为固定常数 $x$，再追问这一行动对某些响应变量 $Y$ 的概率会产生什么效应。本章引入这一分析的若干推广。

首先（\S 4.1），我们讨论行动相对于观测在概率论、决策分析与因果建模中的地位，并提出这样一个论题：因果模型的主要作用，在于便利地评估那些在模型构建时未曾预料到的\emph{新颖}行动与政策的效应。

\S 4.2 把第三章的识别分析推广到“若见到 $z$ 就执行 $x$”形式的条件行动，以及“若见到 $z$ 就以概率 $p$ 执行 $x$”形式的随机策略。我们将看到，这些更精细的干预的评估与识别，都可以由原始干预的分析得到。\S 4.3 用第三章建立的干预演算，给出一类半马尔可夫模型的图形刻画——在这类模型中，一个变量对另一个变量的因果效应可被识别。

\S 4.4 处理序贯计划的效应评估问题，即由若干随时间变化的行动（有些是同时采取的）构成的、旨在产生某一结果的序列。我们给出一种图形方法，用以从非实验研究中估计此类计划的效应——在这类研究中，某些行动受观测与前序行动的影响，某些观测又受行动的影响，而某些混杂变量未被测量。我们将说明，把一个计划分解为其构成行动来分析，比把行动集合当作单一实体处理有实质性的好处。

最后，\S 4.5 讨论如何区分直接效应与间接效应。我们将说明，直接效应可以用 \S 4.4 所发展的图形方法识别。一个关于大学录取中所谓性别歧视的例子，将用以演示正确分析直接效应所需的假设。

\section{导论}

\subsection{行动、行为与概率}

行动有两种解释：反应式的与审思式的。反应式解释把行动视为行动者之信念、性情与环境输入的后果，如“亚当吃了苹果，因为夏娃递给了他”。审思式解释则把行动视为经过权衡的决策中的一个选项，通常涉及对后果的比较，如“亚当在想，倘若他吃了苹果，上帝会怎么做”。我们把前者称为“行为”（act），后者称为“行动”（action），以区分这两种观点。行为是从外部看的，行动是从内部看的。因此，行为可以被预测，也可以作为行动者所受刺激与动机的证据（只要行动者是我们模型的一部分）。相反，行动既不能被预测，也不能提供证据，因为（按定义）它们尚在斟酌之中，一旦执行便转化为行为。

行动与行为的混淆，导致了纽科姆悖论（Newcomb's paradox，Nozick 1969），以及所谓\emph{证据决策论}（evidential decision theory）中的其他种种怪论——这种理论鼓励决策者把某一行动一旦付诸实施将会提供的证据纳入考虑。这一奇特理论似乎源自 Jeffrey 那本颇具影响的《决策的逻辑》（\textit{The Logic of Decision}，Jeffrey 1965）：书中把行动当作普通事件（而非干预）处理，于是行动的效应便通过条件化、而不是通过 $\doop(x)$ 这类修改机制的运算来获得。（见 Stalnaker 1972; Gibbard and Harper 1976; Skyrms 1980; Meek and Glymour 1994; Hitchcock 1996。）

常识决策论\footnote{我特意避开“因果决策论”这一流行名称，以免哪怕一丝暗示：任何替代性的、非因果的理论也可以用来指导决策。}指导理性行动者选择使期望效用最大化的选项 $x$\footnote{按照 Stalnaker (1972) 的建议，Gibbard and Harper (1976) 在 $U(x)$ 中使用的是 $P(x\leadsto y)$，而不是 $P(y\mid\doop(x))$，其中 $x\leadsto y$ 表示虚拟条件句“若为 $x$，则 $y$”。两种算子的语义密切相关（见 \S 7.4），但 $\doop(x)$ 算子的“删去方程”解释歧义更少，且明确压制了由果推因的推断。}：
\[
U(x) = \sum_y P(y\mid\doop(x))\,u(y),
\]
其中 $u(y)$ 是结局 $y$ 的效用；相反，“证据决策论”要求最大化条件期望
\[
U_{\mathrm{ev}}(x) = \sum_y P(y\mid x)\,u(y),
\]
其中 $x$ 被（不恰当地）当作一个被观测到的命题。

这一谬误所产生的悖论显而易见：病人应当避免去看医生，“以降低自己患重病的概率”（Skyrms 1980, p.~130）；工人绝不应当赶着上班，以降低自己睡过了头的概率；学生不应当复习备考，免得这反证自己功课落后；如此等等。一言以蔽之，一切补救行动都应被禁绝，免得它们提高了“确实需要补救”的概率。

这类逻辑的怪诞之处，源于把行动当作受过去关联支配的行为，而不是当作自由选择的对象——而 $\doop(x)$ 算子的语义恰恰要求后者。这种“证据式”决策论说教道：人绝不应忽视真正的统计证据（在我们这里，即一个行为通常提供的、关于该行为是否必要的证据）；但正当的决策论提醒我们：行动——按其定义本身——使这类证据与当前的决策无关，因为行动改变了行为通常所服从的那些概率。\footnote{这类证据在行动者本人的概率空间内确实变得无关；然而在多主体决策情境中，每个主体都理应知晓其他主体会如何解释他那些待决的“将会如此”的行为。}

这段故事的教益，可以用下面这首助记韵语概括：
\begin{center}
\kaishu
无论一个行为可能提供怎样的证据，\\[2pt]
关于究竟是什么引起了该行为，\\[2pt]
都绝不该被用来帮助人决定：\\[2pt]
是否要选择这同一个行为。
\end{center}

证据决策论在哲学文献中只是短暂的一瞬，今天已没有哲学家认真对待该理论的原始版本。尽管如此，近来仍有人试图复活 Jeffrey 的期望效用：用 $P(y\mid x,K)$ 取代 $P(y\mid x)$，其中 $K$ 表示各种背景情境，选取它是为了压制虚假关联（如式 \textup{(3.13)} 那样）（Price 1991; Hitchcock 1996）。这类尝试回响着一种过分狭隘的经验论传统：按此传统，理性行动者只靠一种信息来源——统计关联——而生，也只因之而死，因而期望效用除了贝叶斯条件化之外不应容许任何其他运算。这一传统正迅速让位于一种更宽容的观念：理性行动者应当按照行动的理论行事；自然地，这类理论要求针对行动的条件化（如 $\doop(x)$），而把贝叶斯条件化留给对被动观测的表示（见 Goldszmidt and Pearl 1992; Meek and Glymour 1994; Woodward 1995）。

原则上，行动并不属于概率论的一部分，这也不难理解：概率刻画世界中正常的关系，而行动表示对这些关系加以扰动的干预。因此，在整个概率论与统计学文献中，行动被当作外来实体，也就不足为奇了：它们既不作概率表达式的变元，也不作对这类表达式作条件化的事件。

即便在统计决策论文献中（如 Savage 1954）——行动是分析的主要对象——赋予行动的符号也仅仅充当区分不同概率函数的指标，而不是与概率所定义的那些变量处于逻辑关系之中的实体。Savage (1954, p.~14) 把“行为”定义为“给世界的每一状态附上一个后果的函数”，并把一个决策引向另一个决策的决策链当作单个决策来处理。然而，引导我们从更基本的考量推断行动与策略之后果的那种逻辑，却被排除在形式体系之外。例如，考虑“增税”“减税”与“提高利率”这三个行动。三者的后果都必须在分析之前分别指明；任何一个都不能由其他几个推断出来。于是，若给定两个概率 $P_A$ 与 $P_B$，分别表示在行动 $A$ 与行动 $B$ 之下通行的概率，我们无论如何也无法由这一输入推出对应于联合行动 $A\wedge B$ 的概率 $P_{A\wedge B}$，更不能推出 $A$ 与 $B$ 之任何布尔组合的概率。这意味着，原则上一切预期到的联合行动的影响都必须事先指明——这是一项无法完成的任务。

行动在概率论中的特殊地位，最清楚不过地体现在与观测地位的对比之中。通过在世界的各种可能状态上指定一个概率函数 $P(s)$，我们就自动指明了概率应当如何随每一个可设想的观测 $e$ 而改变，因为 $P(s)$ 使我们能够（通过对 $e$ 取条件）算出每一对事件 $E$ 与 $e$ 的后验概率 $P(E\mid e)$。然而，指定 $P(s)$ 丝毫无助于说明概率应当如何随一个外部行动 $\doop(A)$ 而改变。一般地，若要把行动 $\doop(A)$ 描述为一个把 $P(s)$ 变换为 $P_A(s)$ 的函数，那么 $P(s)$ 对 $P_A(s)$ 的性质一无所告，即便 $A$ 是一个 $P(A)$ 有定义的初等事件（例如“把温度升高 1 度”或“打开洒水器”）。除了“若 $s$ 蕴含 $\neg A$ 则 $P_A(s)$ 为零”这一平凡要求——它对每个 $P(s)$ 一律适用——之外，概率论并不告诉我们 $P_A(s)$ 应当如何不同于 $P_{A'}(s)$，其中 $P'(s)$ 是另一个行动前的概率函数。对 $A$ 取条件显然不足以刻画这一变换，正如我们在第一、三章的许多例子中所见（例如见 \S 1.3.1），因为条件化表示的是不变世界中的被动观测，而行动改变世界。

借助视觉感知的类比，我们可以说，$P(s)$ 所含的信息类似于对一个三维物体的精确描述：它足以预测从该物体之外的任何角度将如何看到它，却不足以预测当该物体被外力摆弄、挤压时将如何被看到。要作出这类预测，必须补充关于该物体物理性质的信息。类似地，描述从 $P(s)$ 到 $P_A(s)$ 之变换所需的补充信息，应当指明世界中那些在行动 $\doop(A)$ 之下保持不变的元素。这种额外信息由因果知识提供，而 $\doop(\cdot)$ 算子使我们能够通过局部修改图或结构方程来捕捉那些不变元素（从而定义 $P_A(s)$）。下一节将把这一装置与标准决策论处理行动的方式作比较。

\subsection{决策分析中的行动}

传统做法不是把新算子引入概率演算，而是把“看”与“做”的差别归因于可获证据总量的差别。试看两个陈述：“观测到气压计读数为 $x$”与“把气压计设定在水平 $x$”。前者帮助我们预测天气，后者则不能。第一个陈述所描述的证据仅限于气压计的读数，而第二个陈述还告诉我们：气压计被某个主体操纵过；对这额外证据取条件，应当使气压计读数与预测降雨变得无关。

这一做法的实际内容，是把行动的主体作为变量纳入分析，构造一个包含这些主体之决策的增广分布函数，再通过对这些决策变量取特定值来推断行动的效应。例如，操纵气压计的主体可以作为决策变量“挤压气压计”进入系统；把这个变量并入概率分布之后，我们只需在增广分布上对事件“气压计以力 $y$ 被挤压并达到水平 $x$”取条件，便可推断操纵气压计的影响。

要使这一条件化方法在评估未来行动的效应时正确运作，必须把操纵主体视为出于自由意志行事的理想实验者，并把相应的决策变量视为外生的——即不受系统中其他变量的因果影响。例如，若增广概率函数编码了这样一个事实：气压计的现任主人每次感到关节炎疼痛时往往就会去挤压气压计，那么我们就无法用该函数来评估有意挤压气压计的效应，即便是同一主人所为。回想行为与行动之别：每当我们着手计算一个待决行动的效应时，必须忽略过去曾约束或触发该行动之执行的一切机制。相应地，事件“气压计被挤压过”必须作为与操纵时刻之前发生的一切事件相互独立者进入增广概率函数，类似于行动变量 $F$ 在图 3.2 中进入增广网络的方式。

这一方案恰好对应于决策分析中处理行动的方式，如影响图（influence diagrams，IDs）文献所描绘的那样（Howard and Matheson 1981; Shachter 1986; Pearl 1988b, chap.~6; Dawid 2002）。每个决策变量被表示为一个外生变量（图中的无父节点），它对其他变量的影响则以条件概率评估并编码，与图中任何其他父节点的影响类似。\footnote{影响图文献坚持把图中的连线与任何因果解释相剥离（Howard and Matheson 1981; Howard 1990），这与流行做法相悖。正是因果解释才使我们能够把决策变量当作根节点，并为分析构造恰当的决策树；演示见 \S 11.6。}

这一做法的困难在于：我们必须事先预料到、并显式表示出将来可能想要评估其效应的一切行动。这使建模过程变得过分繁琐，甚至完全无法管理。以电路诊断为例，把每一种可设想的元件更换行为（同样地，每一种可设想的与电压源、电流源等的连接）都表示为图中的一个节点，将十分笨拙。相反，这类更换的效应本就隐含在电路图本身之中，只要给定其因果解释，便可从图中推出。计量经济建模同样如此：把每一种可设想的政策干预变体都表示为经济方程中的一个新变量，将十分笨拙。相反，只要能把每项政策的即时效应与方程中相应的变量和参数联系起来，这类干预的效应便可从方程的结构解释中推出。例如，复合行动“增税并降低利率”不必作为新变量引入方程，因为只要“税负水平”与“利率”已经作为（外生或内生）变量出现在方程中，该行动的效应便可推出。

无需事先枚举干预便能预测其效应，这是我们从因果建模中获得的主要优势之一，也是因果概念所发挥的主要功能之一。由于行动或行动组合的数量极其庞大，它们无法在模型中显式表示，而只能以每个行动直接强制为真的命题作为索引。强制这些命题为真所带来的间接后果，则由模型中各变量之间的因果关系推断出来。我们将在第七章（\S 7.2.4）回到这一主题，那里将进一步探讨这一编码方案运作所必须满足的不变性假设。

\subsection{行动与反事实}

作为贝叶斯条件化的一种替代，哲学家（Lewis 1976; G\"ardenfors 1988）研究了另一种称为“映像”（imaging）的概率变换。它在分析虚拟条件句时被认为有用，也能更恰当地表示与行动相联系的变换。贝叶斯条件化 $P(s\mid e)$ 把被 $e$ 排除的各状态上的全部概率质量转移到其余状态（按它们当前概率 $P(s)$ 的比例分配）；映像则不同：每个被排除的状态 $s$ 把自身的质量单独转移到一组被认为与 $s$ “最接近”的状态 $S^*(s)$ 上（见 \S 7.4.3）。映像虽为行动提供了更恰当、更一般的框架（Gibbard and Harper 1976），却几乎没有对选择函数 $S^*(s)$ 的精确规定施加任何约束。于是，枚举未来行动的问题，就被替换为如何编码状态之间距离的问题——这种编码既要经济，又要尊重人们对领域中运作之因果律的通常理解。考虑到行动的间接波及往往导致与出发点颇为不同的世界（Fine 1975），后一要求绝非轻而易举。

使“最接近世界”进路符合因果律（\S 7.4）所遇到的困难，在本书所采取的结构进路中被绕开了：我们直接把干预概念建立在因果机制之上，并利用伴随这些机制的不变性与自主性。这种“机制修改”进路可以看作“最接近世界”进路的一个特例，其中相近性度量被精心设计，以尊重领域中的因果机制；由此得到的选择函数 $S^*(s)$ 表示在式 \textup{(3.11)} 中（见其后的讨论）。

这种机制修改语义的可操作性已在第三章得到演示，并带来了对行动效应的定量预测——包括在模型构建时未曾考虑过的行动与行动组合；而在构建模型时，建模者可以自由地描述自然如何运作，不必为外部干预的念头所累。第七章将进一步用机制修改的解释为反事实陈述提供语义，如 \S 1.4.4 所概述。本章则把 do 演算的应用推广到复杂政策的分析与效应的分解。

\section{条件行动与随机策略}

我们在识别分析（\S 3.3—3.4）中所考虑的干预，仅限于迫使一个变量或一组变量 $X$ 取某个指定值 $x$ 的行动。一般地（见 \S 3.2.3 的过程控制例子），干预可能涉及复杂的政策：使变量 $X$ 以特定方式响应另一组变量 $Z$——例如通过函数关系 $x=g(z)$，或通过随机关系，使 $X$ 以概率 $P^*(x\mid z)$ 被置为 $x$。我们将依据 Pearl (1994b) 说明：识别这类政策的效应，等价于识别表达式 $P(y\mid\hat{x},z)$。

以 $P(y\mid\doop(X=g(z)))$ 表示政策 $\doop(X=g(z))$ 之下通行的（$Y$ 的）分布。为计算 $P(y\mid\doop(X=g(z)))$，我们对 $Z$ 取条件，并写出
\begin{equation*}
\begin{aligned}
P(y\mid\doop(X=g(z))) &= \sum_z P(y\mid\doop(X=g(z)),z)\,P(z\mid\doop(X=g(z)))\\
&= \sum_z P(y\mid\hat{x},z)\big|_{x=g(z)}\,P(z)\\
&= E_z\big[P(y\mid\hat{x},z)\big|_{x=g(z)}\big].
\end{aligned}
\end{equation*}
其中等式
\[
P(z\mid\doop(X=g(z))) = P(z)
\]
自然来自这样的事实：$Z$ 不可能是 $X$ 的后代；因此，对 $X$ 施加的任何控制都不会影响 $Z$ 的分布。于是我们看到，政策 $\doop(X=g(z))$ 的因果效应可以直接由 $P(y\mid\hat{x},z)$ 的表达式求得：只需把 $g(z)$ 代入 $x$，并对 $Z$ 取期望（使用观测分布 $P(z)$）。

这一条件政策的可识别性准则，比无条件干预的准则略为严格。显然，若政策 $\doop(X=g(z))$ 可识别，则简单干预 $\doop(X=x)$ 也可识别，因为只要取 $g(z)=x$ 便得到后者。反之则不成立：对 $Z$ 取条件可能造成某些依赖，从而使 $P(y\mid\hat{x},z)$ 无法成功化简为无帽表达式。Kuroki and Miyakawa (1999a, 2003) 给出了图形准则。

随机策略给 $x$ 强加一个新的条件分布 $P^*(x\mid z)$，可用类似方式处理。我们把随机干预视为一个随机过程：其中无条件干预 $\doop(X=x)$ 以概率 $P^*(x\mid z)$ 被施行。于是，给定 $Z=z$，干预 $\doop(X=x)$ 将以概率 $P^*(x\mid z)$ 发生，并产生由 $P(y\mid\hat{x},z)$ 给出的因果效应。对 $x$ 与 $z$ 取平均，便得到随机策略 $P^*(x\mid z)$ 的（对 $Y$ 的）效应：
\begin{equation*}
P(y)\big|_{P^*(x\mid z)} = \sum_x\sum_z P(y\mid\hat{x},z)\,P^*(x\mid z)\,P(z).
\end{equation*}
由于 $P^*(x\mid z)$ 是外部指定的，我们又一次看到：$P(y\mid\hat{x},z)$ 的可识别性，是任何依 $Z$ 的取值来塑造 $X$ 之分布的随机策略可识别的充分必要条件。

在规划中特别重要的是一类 STRIPS 式行动（Fikes and Nilsson 1971）：其即时效果 $X=x$ 取决于某组变量 $W$ 上的某个启用前提 $C(w)$ 是否满足。为表示这类行动，令 $Z=W\cup\PA_X$，并设定
\begin{equation*}
P^*(x\mid z)=
\begin{cases}
P(x\mid\pa_X), & \text{若 } C(w)=\text{false},\\
1, & \text{若 } C(w)=\text{true 且 } X=x,\\
0, & \text{若 } C(w)=\text{true 且 } X\neq x.
\end{cases}
\end{equation*}

\section{行动的效应何时可识别？}

第三章中，我们建立了若干图形准则，用以在存在未观测变量时判定一个变量对另一个变量的效应 $P(y\mid\doop(x))$ 是否可识别。这些准则——如后门准则（定理 3.3.2）与前门准则（定理 3.3.4）——属于一类更一般的半马尔可夫模型的特别情形：对这类模型反复施用 do 演算的推理规则（定理 3.4.1），可将 $P(y\mid\hat{x})$ 化简为无帽表达式，从而使之可识别。本节刻画一类更广的模型，其中的因果效应 $P(y\mid\hat{x})$ 可识别。这一类被 Tian and Pearl (2002a) 在定理 3.7 中所建立的类、以及其后 Shpitser and Pearl (2006b) 给出的完整刻画所包含。此处引入它，是因为它具有直观的吸引力。

\subsection{识别的图形条件}

定理 4.3.1 以四个图形条件的形式刻画一类模型：当 $X$ 与 $Y$ 是图中的单节点时，其中任一条件都足以保证 $P(y\mid\hat{x})$ 可识别。定理 4.3.2 则断言：要使 $P(y\mid\hat{x})$ 在 do 演算中可识别，模型至少须满足这四个条件之一。鉴于 do 演算的完备性，我们得出结论：对于任何与定义 3.2.4 之语义相容的识别方法，这四个条件之一都是必要的。

\begin{theorem}[Galles and Pearl 1995]
设 $X$ 与 $Y$ 表示以图 $G$ 为特征的半马尔可夫模型中的两个单变量。$P(y\mid\hat{x})$ 可识别的一个充分条件是：$G$ 满足下列四个条件之一。
\begin{enumerate}[label=\arabic*.,leftmargin=2.2em]
\item $G$ 中不存在从 $X$ 到 $Y$ 的后门路径；即 $(X\indep Y)_{G_{\bar{X}}}$。
\item $G$ 中不存在从 $X$ 到 $Y$ 的有向路径。
\item 存在一个节点集 $B$，它阻断了从 $X$ 到 $Y$ 的一切后门路径，使得 $P(b\mid\hat{x})$ 可识别。（本条件的一个特例是：$B$ 完全由 $X$ 的非后代构成，此时 $P(b\mid\hat{x})$ 直接化简为 $P(b)$。）
\item 存在节点集 $Z_1$ 与 $Z_2$，使得：
\begin{enumerate}[label=(\roman*),leftmargin=2.6em]
\item $Z_1$ 阻断从 $X$ 到 $Y$ 的每一条有向路径（即 $(Y\indep X\mid Z_1)_{G_{\bar{Z}_1\bar{X}}}$）；
\item $Z_2$ 阻断 $Z_1$ 与 $Y$ 之间的一切后门路径（即 $(Y\indep Z_1\mid Z_2)_{G_{\bar{X}\bar{Z}_1}}$）；
\item $Z_2$ 阻断 $X$ 与 $Z_1$ 之间的一切后门路径（即 $(X\indep Z_1\mid Z_2)_{G_{\bar{X}}}$）；且
\item $Z_2$ 不激活任何从 $X$ 到 $Y$ 的后门路径（即 $(X\indep Y\mid Z_1,Z_2)_{G_{\bar{Z}_1\bar{X}(\bar{Z}_2)}}$）。（只要 (i)—(iii) 成立，且 $Z_2$ 中无一是 $X$ 的后代，本条件即成立。）
\end{enumerate}
（条件 4 的一个特例是：$Z_2=\varnothing$，且不存在从 $X$ 到 $Z_1$ 或从 $Z_1$ 到 $Y$ 的后门路径。）
\end{enumerate}
\end{theorem}

\noindent\textbf{定理 4.3.1 的证明}\quad
\emph{条件 1}。本条件直接由规则 2（见定理 3.4.1）得出。若 $(Y\indep X)_{G_{\bar{X}}}$，则我们可以立即把 $P(y\mid\hat{x})$ 换成 $P(y\mid x)$，故该查询可识别。

\emph{条件 2}。若 $G$ 中不存在从 $X$ 到 $Y$ 的有向路径，则 $(Y\indep X)_{G_{\bar{X}}}$。于是由规则 3 得 $P(y\mid\hat{x})=P(y)$，故该查询可识别。

\emph{条件 3}。若存在节点集 $B$ 阻断从 $X$ 到 $Y$ 的一切后门路径（即 $(Y\indep X\mid B)_{G_{\bar{X}}}$），则可把 $P(y\mid\hat{x})$ 展开为 $\sum_b P(y\mid\hat{x},b)P(b\mid\hat{x})$，并由规则 2 把 $P(y\mid\hat{x},b)$ 改写为 $P(y\mid x,b)$。若查询 $P(b\mid\hat{x})$ 可识别，则原查询也必可识别。例子见图 4.1。

\emph{条件 4}。若存在节点集 $Z_1$ 阻断从 $X$ 到 $Y$ 的一切有向路径，又存在节点集 $Z_2$ 阻断 $G_{\bar{X}}$ 中 $Y$ 与 $Z_1$ 之间的一切后门路径，则展开
\[
P(y\mid\hat{x})=\sum_{z_1,z_2}P(y\mid\hat{x},z_1,z_2)P(z_1,z_2\mid\hat{x}),
\]
并由规则 2 把 $P(y\mid\hat{x},z_1,z_2)$ 改写为 $P(y\mid\hat{x},\hat{z}_1,z_2)$，因为 $Z_1$ 与 $Y$ 之间的一切后门路径都被 $G_{\bar{X}}$ 中的 $Z_2$ 阻断。由规则 3，可把 $P(y\mid\hat{x},\hat{z}_1,z_2)$ 化简为 $P(y\mid\hat{z}_1,z_2)$，因为 $(Y\indep X\mid Z_1,Z_2)_{G_{\bar{Z}_1\bar{X}(\bar{Z}_2)}}$。若 $(Y\indep Z_1\mid Z_2)_{G_{\bar{Z}_1}}$，则可把 $P(y\mid\hat{z}_1,z_2)$ 改写为 $P(y\mid z_1,z_2)$。这一独立性唯一不能成立的情形，是存在一条从 $Y$ 经 $X$ 到 $Z_1$ 的路径，因为 $(Y\indep Z_1\mid Z_2)_{G_{\bar{X}\bar{Z}_1}}$。然而，我们可以通过对 $X$ 取条件并求和来阻断这条路径，从而得到
\[
\sum_{x'}P(y\mid\hat{z}_1,z_2,x')P(x'\mid\hat{z}_1,z_2).
\]
现在可以用规则 2 把 $P(y\mid\hat{z}_1,z_2,x')$ 改写为 $P(y\mid z_1,z_2,x')$。$P(x'\mid\hat{z}_1,z_2)$ 一项可用规则 3 改写为 $P(x'\mid z_2)$，因为 $Z_1$ 是 $X$ 的子节点且图是无环的。于是该查询可改写为
\[
\sum_{z_1,z_2}\sum_{x'}P(y\mid z_1,z_2,x')P(x'\mid z_2)P(z_1,z_2\mid\hat{x}),
\]
并且有 $P(z_1,z_2\mid\hat{x})=P(z_2\mid\hat{x})P(z_1\mid\hat{x},z_2)$。由于 $Z_2$ 由 $X$ 的非后代构成，可用规则 3 把 $P(z_2\mid\hat{x})$ 改写为 $P(z_2)$。又因 $Z_2$ 阻断从 $X$ 到 $Z_1$ 的一切后门路径，可用规则 2 把 $P(z_1\mid\hat{x},z_2)$ 改写为 $P(z_1\mid x,z_2)$。于是整个查询可改写为 $\sum_{z_1,z_2}\sum_{x'}P(y\mid z_1,z_2,x')P(x'\mid z_2)P(z_1\mid x,z_2)P(z_2)$。例子见图 4.2。\hfill$\blacksquare$

\begin{figure}[htbp]
\centering
\begin{tikzpicture}[causal,x=1.3cm,y=1.3cm,every node/.style={transform shape}]
  % ---- panel (a) ----
  \begin{scope}
    \node[vtx] (bia) at (0,1.6) {$B_1$};
    \node[vtx] (xia) at (1.5,1.5) {$X$};
    \node[vtx] (b2a) at (3.0,0.9) {$B_2$};
    \node[vtx] (yia) at (1.5,0.0) {$Y$};
    \foreach \s/\t in {bia/xia,bia/yia,xia/yia,b2a/yia}\draw[ed] (\s)--(\t);
    \draw[bid] (xia) to[bend left=18] (b2a);
    \node[font=\small] at (1.5,-0.65) {(a)};
  \end{scope}
  % ---- panel (b) ----
  \begin{scope}[xshift=5.4cm]
    \node[vtx] (xib) at (0,1.6) {$X$};
    \node[vtx] (zib) at (1.1,0.75) {$Z$};
    \node[vtx] (bib) at (2.3,0.15) {$B$};
    \node[vtx] (yib) at (0,0.0) {$Y$};
    \foreach \s/\t in {xib/yib,xib/zib,zib/bib,bib/yib}\draw[ed] (\s)--(\t);
    \draw[bid] (xib) to[bend left=32] (bib);
    \node[font=\small] at (1.1,-0.65) {(b)};
  \end{scope}
\end{tikzpicture}
\caption{定理 4.3.1 的条件 3。(a) 中集合 $\{B_1,B_2\}$ 阻断了从 $X$ 到 $Y$ 的一切后门路径，且 $P(b_1,b_2\mid\hat{x})=P(b_1,b_2)$。(b) 中节点 $B$ 阻断了从 $X$ 到 $Y$ 的一切后门路径，且 $P(b\mid\hat{x})$ 可依条件 4 识别。}
\label{fig:4.1}
\end{figure}

\begin{theorem}
定理 4.3.1 的四个条件中至少有一个是可识别性的必要条件。也就是说，若定理 4.3.1 的四个条件在图 $G$ 中全都不成立，则不存在任何有限的推理规则序列，能把 $P(y\mid\hat{x})$ 化简为无帽表达式。
\end{theorem}

定理 4.3.2 的证明见 Galles and Pearl (1995)。

\begin{figure}[htbp]
\centering
\begin{tikzpicture}[causal,x=1.3cm,y=1.3cm,every node/.style={transform shape}]
  % ---- (a) ----
  \begin{scope}
    \node[vtx] (x1a) at (0,1.8) {$X$};
    \node[vtx] (z1a) at (0,0.9) {$Z_1$};
    \node[vtx] (y1a) at (0,0.0) {$Y$};
    \draw[ed] (x1a)--(z1a); \draw[ed] (z1a)--(y1a);
    \draw[bid] (x1a) to[bend right=55] (y1a);
    \node[font=\small] at (0,-0.65) {(a)};
  \end{scope}
  % ---- (b) ----
  \begin{scope}[xshift=3.3cm]
    \node[vtx] (x2b) at (0,1.8) {$X$};
    \node[vtx] (z1b) at (0,0.9) {$Z_1$};
    \node[vtx] (y2b) at (0,0.0) {$Y$};
    \node[vtx] (z2b) at (1.5,1.25) {$Z_2$};
    \draw[ed] (x2b)--(z1b); \draw[ed] (z1b)--(y2b); \draw[ed] (z2b)--(z1b);
    \draw[bid] (x2b) to[bend left=18] (z2b);
    \draw[bid] (z2b) to[bend left=18] (y2b);
    \node[font=\small] at (0,-0.65) {(b)};
  \end{scope}
  % ---- (c) ----
  \begin{scope}[xshift=7.0cm]
    \node[vtx] (x3c) at (0,1.8) {$X$};
    \node[vtx] (z1c) at (0,0.9) {$Z_1$};
    \node[vtx] (y3c) at (0,0.0) {$Y$};
    \node[vtx] (z2c) at (1.5,1.25) {$Z_2$};
    \draw[ed] (x3c)--(z1c); \draw[ed] (z1c)--(y3c);
    \draw[ed] (x3c)--(z2c); \draw[ed] (z2c)--(z1c); \draw[ed] (z2c)--(y3c);
    \draw[bid] (x3c) to[bend left=18] (z2c);
    \draw[bid] (z2c) to[bend left=18] (y3c);
    \node[font=\small] at (0,-0.65) {(c)};
  \end{scope}
\end{tikzpicture}
\caption{定理 4.3.1 的条件 4。(a) 中 $Z_1$ 阻断了从 $X$ 到 $Y$ 的一切有向路径，且空集阻断了 $G_{\bar{X}}$ 中从 $Z_1$ 到 $Y$ 的一切后门路径、以及 $G$ 中从 $X$ 到 $Z_1$ 的一切后门路径。(b) 与 (c) 中，$Z_1$ 阻断了从 $X$ 到 $Y$ 的一切有向路径，$Z_2$ 阻断了 $G_{\bar{X}}$ 中从 $Z_1$ 到 $Y$ 的一切后门路径、以及 $G$ 中从 $X$ 到 $Z_1$ 的一切后门路径。}
\label{fig:4.2}
\end{figure}

\subsection{效率方面的评注}

把定理 4.3.1 实现为判定可识别性的系统方法时，条件 3 与条件 4 似乎要求穷举搜索。例如，要证明条件 3 不成立，我们须证明不存在这样的阻断集 $B$。所幸下列定理使我们能大幅剪裁搜索空间，从而使该检验变得可处理。

\begin{theorem}
若 $P(b_i\mid\hat{x})$ 对某个极小集 $B_i$ 可识别，则 $P(b_j\mid\hat{x})$ 对任何其他极小集 $B_j$ 也可识别。
\end{theorem}

定理 4.3.3 使我们能用单个极小阻断集 $B$ 来检验条件 3。若 $B$ 满足条件 3 的要求，则该查询可识别；否则条件 3 无法满足。证明本定理时，我们用到下述引理。

\begin{lemma}
若查询 $P(y\mid\hat{x})$ 可识别，且节点集 $Z$ 位于从 $X$ 到 $Y$ 的某条有向路径上，则查询 $P(z\mid\hat{x})$ 可识别。
\end{lemma}

\begin{theorem}
设 $Y_1$ 与 $Y_2$ 是两个节点子集，使得或者 (i) $Y_1$ 中无一是 $X$ 的后代，或者 (ii) $Y_1$ 与 $Y_2$ 中所有节点都是 $X$ 的后代、且 $Y_1$ 中所有节点都是 $Y_2$ 的非后代。$P(y_1,y_2\mid\hat{x})$ 的化简序列（依推论 3.4.2）存在，当且仅当 $P(y_1\mid\hat{x})$ 与 $P(y_2\mid\hat{x},y_1)$ 二者的化简序列都存在。
\end{theorem}

概率 $P(y_1,y_2\mid\hat{x})$ 若把该程序分别施于 $P(y_2\mid\hat{x},y_1)$ 与 $P(y_1\mid\hat{x})$，或许能通过定理 4.3.1 的检验；但若试图把该检验施于 $P(y_1\mid\hat{x},y_2)$，则找不到规则的化简序列。图 4.3 正是这样的例子。定理 4.3.5 保证：若 $P(y_1,y_2\mid\hat{x})$ 存在化简序列，则只要恰当选择 $Y_1$，我们总能对 $P(y_1\mid\hat{x})$ 与 $P(y_2\mid\hat{x},y_1)$ 二者都找到这样的序列。

\begin{figure}[htbp]
\centering
\begin{tikzpicture}[causal,x=1.3cm,y=1.3cm,every node/.style={transform shape}]
  \node[vtx] (X) at (0,1.9) {$X$};
  \node[vtx] (Y1) at (0,1.0) {$Y_1$};
  \node[vtx] (Z) at (0,0.0) {$Z$};
  \node[vtx] (Y2) at (1.5,0.35) {$Y_2$};
  \draw[ed] (X)--(Y1); \draw[ed] (Y1)--(Z);
  \draw[ed] (Y1)--(Y2); \draw[ed] (Y2)--(Z);
  \draw[bid] (X) to[bend right=55] (Y2);
\end{tikzpicture}
\caption{定理 4.3.1 保证 $P(y_2\mid\hat{x},y_1)$ 与 $P(y_1\mid\hat{x})$ 都有化简序列，而 $P(y_1\mid\hat{x},y_2)$ 却没有。}
\label{fig:4.3}
\end{figure}

\begin{theorem}
若存在集合 $Z_1$ 满足条件 4 中对 $Z_1$ 的一切要求，则由 $X$ 的子节点与 $Y$ 的祖先之交构成的集合，也满足条件 4 中对 $Z_1$ 的一切要求。
\end{theorem}

定理 4.3.6 免去了在定理 4.3.1 的条件 4 中搜索 $Z_1$ 之劳。定理 4.3.3—4.3.6 的证明见 Galles and Pearl (1995)。

\subsection{控制查询的闭式表达式}

定理 4.3.1 所定义的算法不仅判定控制查询的可识别性，还给出 $P(y\mid\hat{x})$ 用被观测概率分布表示的闭式表达式（当这样的闭式存在时），如下。

\begin{techpoint}[ClosedForm 算法]
\textbf{函数}：ClosedForm($P(y\mid\hat{x})$)\\
\textbf{输入}：形如 $P(y\mid\hat{x})$ 的控制查询。\\
\textbf{输出}：或者是一个只用被观测变量表示 $P(y\mid\hat{x})$ 的闭式表达式，或者当查询不可识别时返回 FAIL。
\begin{enumerate}[label=\arabic*.,leftmargin=1.8em]
\item 若 $(X\indep Y)_{G_{\bar{X}}}$，则返回 $P(y)$。
\item 否则，若 $(X\indep Y)_{G_{\underline{X}}}$，则返回 $P(y\mid x)$。
\item 否则，令 $B=\mathrm{BlockingSet}(X,Y)$，$Pb=\mathrm{ClosedForm}(P(b\mid\hat{x}))$；若 $Pb\neq\mathrm{FAIL}$，则返回 $\sum_b P(y\mid b,x)*Pb$。
\item 否则，令 $Z_1=\mathrm{Children}(X)\cap(Y\cup\mathrm{Ancestors}(Y))$，$Z_3=\mathrm{BlockingSet}(X,Z_1)$，$Z_4=\mathrm{BlockingSet}(Z_1,Y)$，且 $Z_2=Z_3\cup Z_4$；若 $Y\notin Z_1$ 且 $X\notin Z_2$，则返回
\[
\sum_{z_1,z_2}\sum_{x'}P(y\mid z_1,z_2,x')P(x'\mid z_2)P(z_1\mid x,z_2)P(z_2).
\]
\item 否则，返回 FAIL。
\end{enumerate}
\end{techpoint}

步骤 3 与 4 调用函数 BlockingSet$(X,Y)$，它选取一个把 $X$ 与 $Y$ d-分离的节点集 $Z$。这样的集合可在多项式时间内找到（Tian et al. 1998）。步骤 3 含有一个对算法 ClosedForm 自身的递归调用，以获得因果效应 $P(b\mid\hat{x})$ 的表达式。

\subsection{小结}

定理 4.3.1 的各条件扩展了可识别模型类（如图 3.8 所描绘者）与不可识别模型类（图 3.9）之间的边界。这些条件导致一个有效算法，用以判定 $P(y\mid\hat{x})$ 这类控制查询的可识别性，其中 $X$ 是单个变量。该算法还给出因果效应 $P(y\mid\hat{x})$ 用可估计概率表示的闭式表达式。

尽管 Shpitser and Pearl (2006a) 的完备性结果如今已精确刻画了可识别模型与不可识别模型之间的边界（见定理 3.6.1 之后的讨论），定理 4.3.2 的条件仍可能因其简明与直观而有用。
